(1) 单体热波长 \(\lambda_1=h/\sqrt{2\pi mk_BT}\)；二聚体质量加倍，\(\lambda_2=\lambda_1/\sqrt2\)。取分离原子能量为零基准，二聚体能量为 \(-\varepsilon\)，所以单粒子配分函数 \(z_1=g_1V/\lambda_1^3\)，\(z_2=(g_2V/\lambda_2^3)e^{\varepsilon/(k_BT)}\)。结合越牢固越偏向形成二聚体，故其 Boltzmann 权重应为正指数，不能把 \(-\varepsilon\) 错写成额外的正内能。

(2) 经典理想混合物各自化学势为 \(\mu_j=k_BT\ln(n_j/z_j)\)。反应平衡要求 \(\mu_2=2\mu_1\)，于是 \(n_2/z_2=(n_1/z_1)^2\)，即 \(n_2/V=K_c(T)(n_1/V)^2\)。平衡常数 \(K_c=(g_2/g_1^2)(\lambda_1^6/\lambda_2^3)e^{\varepsilon/(k_BT)}=(g_2/g_1^2)2^{3/2}\lambda_1^3e^{\varepsilon/(k_BT)}\)。它的量纲是体积，因右边浓度平方乘 \(K_c\) 必须仍为浓度；混用无量纲标准态平衡常数会导致量纲错误。

(3) 令 \(n_2=K_cn_1^2/V\)，核数守恒变为 \(2(K_c/V)n_1^2+n_1-N=0\)。唯一非负根是 \(n_1=[V/(4K_c)](\sqrt{1+8K_cN/V}-1)\)，适宜数值计算时也可有理化成 \(n_1=2N/[1+\sqrt{1+8K_cN/V}]\)。再得 \(n_2=(N-n_1)/2\)。由于 \(K_c\propto T^{-3/2}\exp(\varepsilon/k_BT)\)，高温或低密度时 \(K_cN/V\ll1\) 导致 \(n_1\to N\)，低温且模型仍满足经典假设时 \(K_cN/V\gg1\) 导致 \(n_1\ll N\)、几乎全为二聚体。

(4) 两种独立理想平动组分贡献总压强 \(P=(n_1+n_2)k_BT/V=[(N+n_1)/2]k_BT/V\)，完全二聚时压强约为未反应原子气体的一半，体现的是粒子数变化。忽略内部激发后平均能量为 \(U=\frac32(n_1+n_2)k_BT-\varepsilon n_2\)；第二项是负结合能，随 \(T\) 改变的化学组成会向热容加入反应热效应。求实际 \(C_V=dU/dT\) 时必须包含 \(dn_j/dT\)，而不能把 \(n_1,n_2\) 当成彼此独立的常数；模型在低温进入简并或冻结反应动力学时会失效。
